% ============================================================================
%  Part V: How to Implement
% ============================================================================
\section{How to Implement}
\label{sec:how}

This part walks through the practical decisions a generator implementation must make. The reference implementations in Part \ref{sec:implementations} embody these decisions; this section explains why each choice was made so a porter can confidently produce a conformant implementation in any language.

\subsection{The CSPRNG is the bottleneck}

Across every reference implementation, the dominant cost is the call into the operating system's cryptographically-strong PRNG. On Linux this is \code{getrandom(2)}; on macOS and the BSDs, \code{arc4random\_buf(3)} or equivalent; on Windows, \code{BCryptGenRandom} or \code{RtlGenRandom}. The typical cost of one call --- regardless of how few or how many bytes you fetch --- is a few hundred nanoseconds, dominated by the syscall overhead and the kernel-side entropy mixing.

Two practical consequences:

\textbf{Fetch in bulk.} A naive implementation calls the CSPRNG twice per generation: once for the counter step byte and once for the 32-byte random tail. The reference implementations fetch all 33 bytes in a single call (\code{rnd[0]} for the counter step, \code{rnd[1..32]} for the tail), saving one syscall per generation. This is a measurable win.

\textbf{Don't call the CSPRNG from a hot loop.} If you need to generate many TesseraIDs in succession (e.g., backfilling a large table), batch the CSPRNG calls. The reference Python ``fast'' variant pre-fills a buffer and serves IDs out of it. The Rust and Go implementations don't bother because their syscall overhead is already low.

\subsection{Clock clamping}

If the wall clock moves backward (NTP \citep{rfc5905} correction, leap second, system clock adjustment), we use the last-known timestamp instead. This preserves local monotonicity at the cost of generating IDs whose timestamps are slightly ahead of the actual wall clock for as long as the regression lasts. This is the right tradeoff: a non-monotonic ID stream breaks B-tree indexes, replication ordering, and human intuition; a brief future-skew of a few hundred milliseconds is invisible to virtually all consumers.

A more sophisticated implementation could use a monotonic clock to detect regressions and recompute timestamps based on the monotonic offset. The reference implementations don't do this because the simple ``clamp to last\_ts'' approach is correct, simple, and never visible in practice except during clock-correction events.

\subsection{Counter exhaustion}

The 16-bit counter holds 65{,}536 values. With a random step in $[1, 16]$, the average counter value after $N$ generations within a millisecond is $N \cdot 8.5$. This means the counter exhausts at approximately $7{,}700$ generations per millisecond per generator, i.e., about $7.7$ million IDs per second per generator from a single thread.

In practice, no single-threaded generator hits this rate. The CSPRNG syscall caps single-threaded throughput well below 7M ops/sec in every environment measured to date (current per-language figures live in the \texttt{benchmarks/} directory). The exhaustion case is therefore largely theoretical for single-threaded use.

When it does happen --- typically because the system clock is frozen or moving backward and the generator is being driven hard --- the implementation spins in a tight loop waiting for \code{now\_ms()} to advance, then reseeds the counter from \code{rnd[0]}. The spin is bounded because clocks are guaranteed to advance eventually; in practice the spin lasts at most one millisecond.

\subsection{Thread safety}

The Standard generator is lock-free per-instance: it maintains private \code{last\_ts} and \code{counter} state and is safe to call from a single thread. For multi-threaded use, the reference implementations adopt one of two strategies:

\begin{itemize}
  \item \textbf{One generator per thread} (Go's preferred approach): cheap, no contention, and the per-thread state is small enough that even thousands of threads fit comfortably in memory.
  \item \textbf{One global generator behind a mutex} (Python and TypeScript's default): simpler API at the cost of contention. The contention is usually tolerable because the critical section is microseconds.
\end{itemize}

The HLC generator must serialize access in all cases because the HLC algorithm reads-then-writes the \code{(l, c)} state. The reference C++ implementation guards the HLC state with a \code{std::mutex}; the Python implementation uses a \code{threading.Lock}; the Go implementation uses a \code{sync.Mutex}. The contention cost is similar to the Standard global-generator approach.

\subsection{Encoding implementation}

The Base64URL encoding is straightforward: standard 6-bits-per-char packing with the URL-safe alphabet, no padding. The trailing two bytes (offsets 39--40) become three characters carrying 16 bits of value plus 2 bits of zero padding, and decoders MUST verify that the padding bits are zero.

The Crockford Base32 encoding is similar but uses 5-bits-per-char packing and the Crockford alphabet (\code{0--9} then \code{A--Z} minus \code{I/L/O/U}). The decoder MUST be case-insensitive and MUST canonicalize the I/L/O/0 confusables: \code{I} and \code{L} (both cases) decode to value 1; \code{O} (both cases) decodes to value 0. The trailing byte (offset 40) becomes two characters carrying 8 bits of value plus 2 bits of zero padding.

Both encodings are pure functions of the input bytes; no random state, no time. They can be implemented in roughly 50 lines of code per direction in any modern language.

\subsection{Decoder validation}

A conformant decoder MUST:

\begin{itemize}
  \item Reject inputs whose length is wrong (55 chars for Base64URL, 66 for Crockford, 41 bytes for raw binary).
  \item Reject inputs containing characters outside the alphabet.
  \item Reject inputs whose final character carries non-zero padding bits.
  \item For binary inputs, validate that the variant byte is \code{0x01} or \code{0x02}; reject other values.
\end{itemize}

A conformant decoder SHOULD also expose the variant byte, timestamp, and counter as accessors so that consumers can inspect the structured fields without re-parsing the entire identifier.

